\chapter{El teorema de Tales}\label{ch:g9:thales}

¿Cómo se mide la altura de un \omterm{def:g8:solids:def}{pirámide} sin subirlo? Tales
Mileto comparó su sombra con la sombra de un palo. Su teorema
--- \omterm{def:g6:lines:perp}{lineas paralelas} cortar \omterm{def:g6:lines:objects}{segmentos} en \omterm{def:g9:linfunc:linear}{proporcional} piezas --- es el
corazón matemático de \omterm{def:g7:prop:scale}{escala} maquetas, mapas y ampliaciones, y uno de
Los teoremas más antiguos con nombre.

\section{El teorema}

\begin{theorem}[Tales]\label{thm:g9:thales:direct}
Sean dos rectas que se cruzan en un punto $A $. Llevar $ M $ en la primera línea y $ N $
en el segundo, con $ B $ más adelante en la primera y $ C $ más adelante
el segundo. si las lineas $(MN)$ y $(a.\,C.)$ son \emph{\omterm{def:g6:lines:perp}{paralelo}}, entonces el
triangulos $ AMN $ y $ ABC $ tener \omterm{def:g9:linfunc:linear}{proporcional} lados:
\[
\frac{AM}{AB} = \frac{AN}{AC} = \frac{MN}{a.\,C.}.
\]
\end{theorem}
\admitted

\begin{omfigure}
\begin{tikzpicture}[scale=1.05]
  \coordinate (A) at (0,0);
  \coordinate (B) at (5.0,0.6);
  \coordinate (C) at (3.2,2.8);
  \coordinate (M) at ($(A)!0.55!(B)$);
  \coordinate (N) at ($(A)!0.55!(C)$);
  \draw[thick] (A) -- (B) -- (C) -- cycle;
  \draw[omThm, very thick] (M) -- (N);
  \draw[omThm, very thick] (B) -- (C);
  \fill (A) circle (1.5pt) node[left, font=\small] {$ A $};
  \fill (B) circle (1.5pt) node[right, font=\small] {$ B $};
  \fill (C) circle (1.5pt) node[above, font=\small] {$ C $};
  \fill (M) circle (1.5pt) node[below, font=\small] {$ M $};
  \fill (N) circle (1.5pt) node[above left, font=\small] {$ N $};
\end{tikzpicture}\hspace{1.1cm}%
\begin{tikzpicture}[scale=1.05]
  \coordinate (A) at (0,0);
  \coordinate (B) at (2.6,-1.3);
  \coordinate (C) at (1.4,-2.2);
  \coordinate (M) at ($(A)!-0.62!(B)$);
  \coordinate (N) at ($(A)!-0.62!(C)$);
  \draw[thick] (M) -- (B);
  \draw[thick] (N) -- (C);
  \draw[omThm, very thick] (M) -- (N);
  \draw[omThm, very thick] (B) -- (C);
  \fill (A) circle (1.5pt) node[right=2pt, font=\small] {$ A $};
  \fill (B) circle (1.5pt) node[right, font=\small] {$ B $};
  \fill (C) circle (1.5pt) node[below, font=\small] {$ C $};
  \fill (M) circle (1.5pt) node[left, font=\small] {$ M $};
  \fill (N) circle (1.5pt) node[above, font=\small] {$ N $};
\end{tikzpicture}

{\small Las dos configuraciones de Tales: triángulos anidados (izquierda) y el
``mariposa'', donde $ M $ y $ N $ sentarse al otro lado de $ A $ (bien).
En ambos, $(MN) \parallel (a.\,C.)$ y las tres razones son iguales.}
\end{omfigure}

\begin{method}[Calcular una longitud con Thales]\label{met:g9:thales:compute}
\begin{enumerate}
  \item Identificar el punto $ A $ donde las dos líneas se cruzan y los dos
        \omterm{def:g6:lines:perp}{lineas paralelas}; nombra los dos triángulos;
  \item escribe las tres razones iguales, siempre ``triángulo pequeño sobre grande
        triángulo'';
  \item mantener la igualdad de las dos razones que involucran a los tres conocidos
        longitudes y la desconocida;
  \item resuelve la proporción resultante mediante multiplicación cruzada.
\end{enumerate}
\end{method}

\begin{example}\label{ex:g9:thales:compute}
En la configuración izquierda arriba, supongamos $ AM = 3$, $ AB = 5$, $ AN = 2.4$
y $ MN = 3.3$, y calculemos $ AC $ y $ a.\,C. $. Tales da
\[
\frac{AM}{AB} = \frac{AN}{AC} = \frac{MN}{a.\,C.},
\qquad\text{i.e.}\qquad
\frac35 = \frac{2.4}{AC} = \frac{3.3}{a.\,C.}.
\]
De $\frac35 = \frac{2.4}{AC}$, multiplicando de forma cruzada:
$3 \times AC = 5 \times 2.4 = 12$, entonces $ AC = 4$. De
$\frac35 = \frac{3.3}{a.\,C.}$: $3 \times a.\,C. = 16.5$, entonces $ a.\,C. = 5.5$.
\end{example}

\begin{example}[Tales en la vida real]\label{ex:g9:thales:stick}
Un palo vertical de altura. $1$ m proyecta una sombra de $1.5$ m, mientras que un árbol
proyecta una sombra de $12$ m al mismo tiempo. Los rayos del sol son \omterm{def:g6:lines:perp}{paralelo}, entonces
los triángulos del palo y la sombra y del árbol y la sombra están en Tales
configuración:
\[
\frac{\text{height of tree}}{1} = \frac{12}{1.5},
\qquad\text{so the tree is } 8 \text{ m tall.}
\]
\end{example}

\section{lo contrario}

\begin{theorem}[Converso de Tales]\label{thm:g9:thales:converse}
Con los puntos colocados como en \cref{thm:g9:thales:direct} ($ M $ en $(AB)$,
$ N $ en $(AC)$, en el mismo orden en ambas líneas):
si
\[
\frac{AM}{AB} = \frac{AN}{AC},
\]
luego las lineas $(MN)$ y $(a.\,C.)$ son \omterm{def:g6:lines:perp}{paralelo}.
\end{theorem}
\admitted

\begin{method}[Demostrar o refutar el paralelismo]\label{met:g9:thales:converse}
\begin{enumerate}
  \item Calcule por separado las dos razones. $\dfrac{AM}{AB}$ y
        $\dfrac{AN}{AC}$ (como \omterm{def:g9:fractions:fraction}{fracciones}, no \omterm{def:g6:decimals:rounding}{redondeos});
  \item si son iguales \emph{y} los puntos están en el mismo orden en
        las dos líneas, las líneas son \omterm{def:g6:lines:perp}{paralelo} (conversación de Tales);
  \item si difieren, las líneas son \emph{no} \omterm{def:g6:lines:perp}{paralelo} --- porque si
        Si lo fueran, el teorema de Tales obligaría a que las proporciones fueran iguales.
\end{enumerate}
\end{method}

\begin{example}\label{ex:g9:thales:converse}
En la línea 1, $ AM = 4$ y $ AB = 10$; en la línea 2, $ AN = 6$ y $ AC = 15$.
Entonces
\[
\frac{AM}{AB} = \frac{4}{10} = \frac25,
\qquad
\frac{AN}{AC} = \frac{6}{15} = \frac25 :
\]
proporciones iguales, mismo orden: $(MN) \parallel (a.\,C.)$.

Con $ AC = 14$ en cambio: $\frac{6}{14} = \frac37 \neq \frac25$, entonces el
las líneas no son \omterm{def:g6:lines:perp}{paralelo}.
\end{example}

\section{Ampliación y reducción}

\begin{definition}[Escalar una figura]\label{def:g9:thales:scaling}
Ampliar o reducir una figura en el \emph{factor de escalada}\index{factor de escalada} $ k > 0$ significa multiplicar \emph{todo} es
longitudes por $ k $: un \emph{ampliación} cuando $ k > 1$, a \emph{reducción}
cuando $ k < 1$. En la configuración de Tales, el triángulo $ AMN $ es el
reducción de $ ABC $ por el factor $ k = \frac{AM}{AB}$.
\end{definition}

\begin{proposition}[Efecto sobre ángulos y áreas.]\label{prop:g9:thales:scaling}
Escalado por el factor $ k $ conserva los ángulos y la forma de las figuras,
multiplica cada longitud por $ k $, y multiplica cada \emph{\omterm{def:g6:measure:area}{área}} por
$ k^2$.
\end{proposition}
\admitted

\begin{omfigure}
\begin{tikzpicture}[scale=0.72]
  \draw[thick, omDef, fill=omDef!8] (0,0) rectangle (2,1.2);
  \draw[thick, omThm, fill=omThm!8] (3.4,0) rectangle (7.4,2.4);
  \node[font=\small, omDef] at (1,-0.45) {$2 \times 1.2$};
  \node[font=\small, omThm] at (5.4,-0.45) {$4 \times 2.4$};
  \node[font=\small] at (1,0.6) {area $2.4$};
  \node[font=\small] at (5.4,1.2) {area $9.6$};
\end{tikzpicture}

{\small Duplicar las longitudes ($ k = 2$) multiplica el \omterm{def:g6:measure:area}{área} por
$ k^2 = 4$: cuatro copias del rectángulo pequeño en mosaico del grande.}
\end{omfigure}

\begin{example}
una foto de $10 \times 15$ cm se amplía con \omterm{def:g9:thales:scaling}{factor de escala} $ k = 3$: el
imprimir medidas $30 \times 45$ centímetro. Es \omterm{def:g6:measure:area}{área} va desde $150$ centímetro $^2$ a
$150 \times 9 = 1350$ centímetro $^2$ --- $9$ veces más grande, no $3$.
\end{example}

\section{Ejercicios}

\begin{exercise}[$\star $]\label{exo:g9:thales:1}
las lineas $(MN)$ y $(a.\,C.)$ son \omterm{def:g6:lines:perp}{paralelo}, con $ M $ en $[AB]$ y $ N $ en
$[AC]$; $ AM = 2$, $ AB = 6$, $ AN = 3$, $ a.\,C. = 9$. Calcular $ AC $ y $ MN $.
\end{exercise}

\begin{exercise}[$\star $]\label{exo:g9:thales:2}
Misma configuración: $ AM = 5$, $ MB = 3$ (cuidadoso: $ MB $, no $ AB $!),
$ AN = 4$. Calcular $ AC $.
\end{exercise}

\begin{exercise}[$\star $]\label{exo:g9:thales:3}
En una configuración de mariposa, $ A $ es
entre $ M $ y $ B $ y entre $ N $
y $ C $, con $(MN) \parallel (a.\,C.)$; $ AM = 3$, $ AB = 7.5$, $ AN = 2$,
$ MN = 2.6$. Calcular $ AC $ y $ a.\,C. $.
\end{exercise}

\begin{exercise}[$\star $]\label{exo:g9:thales:4}
$ M $ es
        en $[AB]$ con $ AM = 6$, $ AB = 8$; $ N $ es
        en $[AC]$ con $ AN = 9$,
$ AC = 12$. son las lineas $(MN)$ y $(a.\,C.)$ \omterm{def:g6:lines:perp}{paralelo}?
\end{exercise}

\begin{exercise}[$\star\star $]\label{exo:g9:thales:5}
Lo mismo que arriba con $ AM = 4$, $ AB = 6$, $ AN = 5$, $ AC = 8$. Son $(MN)$
y $(a.\,C.)$ \omterm{def:g6:lines:perp}{paralelo}? Justifique cuidadosamente.
\end{exercise}

\begin{exercise}[$\star\star $]\label{exo:g9:thales:6}
A $1.8$ m persona alta se encuentra $2$ m de distancia de la base de una farola;
sus medidas de sombra $3$ metro. Dibuja la configuración de Tales formada por
lámpara, la persona y la punta de la sombra, y calcular la altura de
la lámpara.
\end{exercise}

\begin{exercise}[$\star\star $]\label{exo:g9:thales:7}
un mapa tiene \omterm{def:g7:prop:scale}{escala} $1 : 25\,000$ (las longitudes en el mapa son las longitudes reales
multiplicado por $ k = \frac{1}{25000}$).
\begin{enumerate}
  \item Dos pueblos son $6.8$ cm de distancia en el mapa. ¿Qué es lo real?
        distancia, en km?
  \item Un bosque tiene un \omterm{def:g6:measure:area}{área} de $8$ centímetro $^2$ en el mapa. cual es su verdadera
        \omterm{def:g6:measure:area}{área}, en kilómetros $^2$?
\end{enumerate}
\end{exercise}

\begin{exercise}[$\star\star $]\label{exo:g9:thales:8}
Triángulo $ ABC $ tiene $ AB = 12$, $ AC = 15$, $ a.\,C. = 18$. el punto $ M $ en
$[AB]$ satisface $ AM = 8$, y la línea que pasa $ M $ \omterm{def:g6:lines:perp}{paralelo} a $(a.\,C.)$
cortes $[AC]$ en $ N $.
\begin{enumerate}
  \item Calcular $ AN $ y $ MN $.
  \item cual es el \omterm{def:g9:thales:scaling}{factor de escala} de $ ABC $ a $ AMN $? comparar el
        \omterm{def:g6:measure:perimeter}{perímetros} de los dos triángulos.
\end{enumerate}
\end{exercise}

\begin{exercise}[$\star\star\star $]\label{exo:g9:thales:9}
Sea $ ABCD $ un trapezoide con $(AB) \parallel (CD)$, $ AB = 4$ y
$ CD = 6$, cuyas diagonales $[AC]$ y $[BD]$ reunirse en $ O $.
Usando Thales en la configuración de mariposa alrededor $ O $, calcular el
relación $\dfrac{OA}{OC}$, y demostrar que $ O $ corta ambas diagonales en el
misma proporción.
\end{exercise}

\section{Problema: La regla sin marcar, la cámara estenopeica y el
tres medios de un trapezoide}

\begin{problem}[{Problema de fin de semana --- Tales en el trabajo: dividir cualquier segmento en partes iguales, medir el Sol con una caja de zapatos y un trapezoide donde se encuentran los tres famosos medios}]
\label{pb:g9:thales:1}
El teorema de Tales es la matemática de \emph{rayos}: rayos de sol,
rayos de luz a través de un agujero de alfiler, rayos de lápiz de un \omterm{def:g5:solids:def}{vértice}. esto
El problema lo utiliza de tres maneras: como herramienta de construcción (dividiendo un
\omterm{def:g6:lines:objects}{segmento} de \emph{cualquier} longitud, \omterm{def:g3:numbers:evenodd}{incluso} $\sqrt2$, en perfectamente iguales
partes), como instrumento de medición (el \omterm{def:g4:geometry:circle}{diámetro} del Sol, con
una caja de zapatos y una moneda), y como lupa en el trapezoide
de \cref{exo:g9:thales:9}, donde la aritmética, geométrica y
Los medios armónicos de esta serie aparecen todos en una sola figura.

\medskip
\noindent\textbf{Parte I --- Dividing with an unmarked ruler.}
\begin{enumerate}
  \item dibuja un \omterm{def:g6:lines:objects}{segmento} $[AB]$ de $7$ centímetro. Para cortarlo en tres
        partes iguales con compás y regla sin marcar: dibuja cualquier rayo
        de $ A $ (no a través de $ B $); bajarse de tres compases iguales
        longitudes en él, dando puntos $ P_1$, $ P_2$, $ P_3$; unirse
        $ P_3$ a $ B $; dibujar el \omterm{def:g6:lines:perp}{paralelas} a $(P_3 B)$ a través de
        $ P_1$ y $ P_2$. Realizar la construcción y marcar donde
        el \omterm{def:g6:lines:perp}{paralelas} cortar $[AB]$.
  \item Justifica con \cref{thm:g9:thales:direct} que los dos
        puntos marcados cortados $[AB]$ exactamente en sus terceros puntos.
  \item Adaptar el método para dividir un \omterm{def:g6:lines:objects}{segmento} en $5$ igual
        partes, luego construir $\frac35$ de un dado \omterm{def:g6:lines:objects}{segmento}.
  \item construir el punto $ M $ de $[AB]$ con
        $\frac{AM}{MB} = \frac23$ (eso es, $ M $ en $\frac25$ de
        el camino desde $ A $). ¿Cuántos pasos iguales en el rayo, y
        que punto se une $ B $?
  \item ¿Por qué esta construcción es mejor que medir con un
        regla graduada? Considere un \omterm{def:g6:lines:objects}{segmento} de longitud $\sqrt2$
        (la diagonal de un cuadrado unitario, irracional por
        \cref{pb:g9:sqrt:1}): qué daría la medición, y
        ¿Qué da Tales?
\end{enumerate}

\medskip
\noindent\textbf{Parte II --- The pinhole camera.}
Introduce un alfiler en uno. \omterm{def:g5:solids:def}{rostro} de un cerrado \omterm{def:g5:solids:def}{caja}: al contrario
\omterm{def:g5:solids:def}{rostro}, aparece una imagen del mundo al revés. un punto de un
objeto, el orificio y el punto de la imagen están alineados --- luz
viaja recto --- entonces el objeto (altura $ H $, a distancia $ D $
delante del agujero) y su imagen (altura $ h $, en la espalda
pared en profundidad $ d $ detrás del agujero) siéntate en la mariposa
configuración de Tales.
\begin{enumerate}[resume]
  \item Deducir la fórmula estenopeica
        $ h = H \times \frac dD $ de Tales en la mariposa
        alrededor del agujero.
  \item A $12$ m soportes de árboles $30$ m de una caja de zapatos de profundidad
        $20$ centímetro. ¿Qué altura tiene su imagen y hacia dónde se eleva?
  \item Señale el \omterm{def:g5:solids:def}{caja} al sol: a $ d = 1$ estoy detrás del
        agujero de alfiler, la imagen del Sol es un disco de \omterm{def:g4:geometry:circle}{diámetro} acerca de
        $9$ mm. Dada la distancia del Sol
        $ D = 1.5 \times 10^{11}$ m, calcule su \omterm{def:g4:geometry:circle}{diámetro} en
        \omterm{def:g9:fractions:scientific}{notación científica} (\cref{pb:g9:fractions:1}'s
        notación en el trabajo). Comparar con el valor real,
        $1.39 \times 10^9$ metro.
  \item La Luna: \omterm{def:g4:geometry:circle}{diámetro} $3.5 \times 10^6$ m, distancia
        $3.8 \times 10^8$ metro. Calcular la razón
        $\frac{\text{diameter}}{\text{distance}}$ para la luna
        y para el sol. ¿Qué afortunada coincidencia tienen los dos?
        Los números revelan --- ¿y qué evento espectacular lo revela?
        hacer posible?
  \item En una o dos frases: ¿por qué aparece la imagen estenopeica?
        boca abajo, y ¿por qué al agrandar el orificio se hace que el
        ¿Imagen más brillante pero más borrosa?
\end{enumerate}

\medskip
\noindent\textbf{Parte III --- The trapezoid of the three
means.}
$ ABCD $ es el trapezoide de \cref{exo:g9:thales:9}:
$(AB) \parallel (CD)$, $ AB = 4$, $ CD = 6$, diagonales que se cruzan en
$ O $ con $\frac{OA}{OC} = \frac{OB}{OD} = \frac{4}{6} =
\frac23$.
\begin{enumerate}[resume]
  \item Práctica de mariposa: dos líneas se cruzan en $ O $ con
        $(MN) \parallel (PQ)$ en posición mariposa; $ OM = 3$,
        $ OP = 7.5$, $ MN = 4$ y $ OQ = 6$. Calcular $ PQ $ y
        $ ON $.
  \item Explica cómo Tales, aplicado con razón. $\frac12$,
        contiene el teorema del punto medio de
        \cref{thm:g8:midpoints:direct} como su caso especial.
  \item Dibuja el \omterm{def:g6:lines:perp}{paralelo} a las bases a través $ O $; se encuentra
        $[d.\,C.]$ en $ E $. Trabajando en el triángulo $ ACD $ (nota
        $\frac{AO}{AC} = \frac25$), calcular $ EO $.
  \item Por el cálculo del espejo en el triángulo $ BDC $, el
        pieza $ OF $ (con $ F $ en $[a.\,C.]$) tiene la misma longitud.
        Concluir que
        \[
        EF = 2 \times \frac{4 \times 6}{4 + 6} = 4.8 :
        \]
        el \omterm{def:g6:lines:perp}{paralelo} a través del cruce diagonal se mide la
        \emph{\omterm{pb:g8:speed:1}{armónico
significar}} de las bases --- el ida y vuelta
        media de \cref{pb:g8:speed:1}, reapareciendo en estado puro
        geometría.
  \item El gran final: calcular las tres medias clásicas de
        las bases $4$ y $6$ --- aritmética
        $\frac{4+6}{2}$, geométrico $\sqrt{4 \times 6}$, armónico
        $\frac{2 \times 4 \times 6}{4 + 6}$ --- y ordenarlos.
        Cada uno vive en el trapezoide: la media aritmética es la
        línea media que une los puntos medios de las piernas
        (\cref{pb:g7:areas:1}), el \omterm{pb:g8:speed:1}{armónico
significar} es tu
        \omterm{def:g6:lines:objects}{segmento} $ EF$, y el \omterm{pb:g8:cosine:1}{geométrico
        significar} es el \omterm{def:g6:lines:perp}{paralelo}
        que corta el trapezoide en dos \emph{similar}
        trapecios (verifique su valor con una calculadora; el
        la prueba es un desafío adicional encantador). ¿Dónde tiene el
        La desigualdad entre los tres ha sido demostrada en este libro.
        (\cref{pb:g8:cosine:1}, \cref{pb:g8:speed:1})?
\end{enumerate}
\end{problem}
